\chapter{Software Design}

\section{Introduction to Software Design}
This chapter turns the requirements of Chapter~3 into a design. It describes the methodology and
process model, the architecture, the behavioral models (activity, sequence and state diagrams), the
static structure (class diagram), the data flows, and the database design with its data dictionary.

\section{Design Methodology and Software Process Model}
\subsection{Design Methodology}
An object-oriented design methodology was used. Domain concepts are modeled as classes (Django
models) with clear responsibilities, and behavior that crosses several classes is placed in service
modules: \code{documents/access.py} (access rules), \code{documents/pipeline/} (processing steps),
\code{assistant/router.py}, \code{assistant/retrieval.py} and \code{assistant/engine.py} (question
answering), \code{audit/service.py} and \code{notifications/service.py}. Every call to an AI model
goes through the \code{ai} package, which hides the provider behind one interface.

Two design principles guided the work:
\begin{enumerate}
\item \textbf{Permission before retrieval.} The function \code{visible\_documents(user)} is the only
way any module obtains documents. Search, lookups, file downloads and answers all start from its result,
so the language model can never receive text from a document the user may not open.
\item \textbf{Evidence for every fact.} Extracted values store their page and verbatim quote, answers
store their sources, and the interface always shows them. A value that cannot be proved is not trusted
automatically.
\end{enumerate}

\subsection{Software Process Model}
The project followed an iterative, incremental Agile process. Work was divided into short iterations,
each ending with a working, tested increment:
\begin{enumerate}
\item Domain re-scoping, data model, RBAC and the four history types.
\item Processing pipeline: extraction, OCR, handwriting, classification, chunking, embeddings.
\item Record field extraction with grounding, the review queue and the records register.
\item Query router, hybrid retrieval and the three answer engines with citations.
\item User interface redesign with the evidence-linked document viewer.
\item Clause classifier fine-tuning, the evaluation suite, reminders and deployment.
\end{enumerate}
This model suited the project because requirements from the supervisor were refined during
development and each increment could be demonstrated.

\section{System Overview}
Figure~\ref{fig:context} shows \sysname{} as a single process with its external entities. Members
(Admins, Managers, Employees and Auditors) provide documents, questions, verifications and
administrative data and receive answers, documents, review items and reminders. The system sends
prompts and page images to the LLM/VLM provider and optionally sends reminder e-mails.

\tikzfig[0.5]{context}{Context Diagram of \sysname{}}{fig:context}

\section{Architectural Design}
Figure~\ref{fig:architecture} shows the architecture. It is a layered client-server architecture with a
separate background worker.

\tikzfig[0.62]{architecture}{System Architecture Diagram of \sysname{}}{fig:architecture}

\begin{enumerate}
\item \textbf{Presentation tier (Vue.js 3 SPA).} Renders all screens, keeps the session token, and
calls the REST API. It contains no business rules; buttons a user may not use are hidden, but the server
enforces every rule again.
\item \textbf{Application tier (Django REST API).} Organized as five Django applications:
\code{accounts} (organizations, members, roles), \code{documents} (documents, versions, pages, chunks,
records, review), \code{assistant} (router, engines, retrieval, conversations), \code{audit} and
\code{notifications}. The \code{ai} package contains the LLM gateway (with deadline and provider
fallback), the embedding model, OCR and the clause classifier.
\item \textbf{Background worker.} A long-running \code{run\_worker} command claims queued jobs from a
database table using \code{SELECT \dots{} FOR UPDATE SKIP LOCKED}, runs the pipeline, retries failures
up to three times and runs the reminder check every hour. Using the database as the queue avoids an extra
message broker and works on Windows and Linux.
\item \textbf{Data tier.} PostgreSQL 16 holds all relational data. The pgvector extension stores chunk
embeddings with an HNSW index (cosine distance), and a weighted \code{tsvector} column with a GIN index
supports full-text search, so permission filters, vector search and keyword search run in the same
database. Uploaded files are stored on disk under organization- and document-specific paths and are
served only through a permission-checked endpoint.
\item \textbf{External services.} An OpenAI-compatible LLM endpoint (Groq by default, Gemini as
fallback) and a vision-language model for handwriting.
\end{enumerate}

\section{Design Models}
\subsection{Activity Diagrams}
\subsubsection{Module 2: Document Ingestion}
Figure~\ref{fig:act-ingest} shows how an uploaded file becomes searchable. Each page chooses its own
reading method, so a PDF that mixes typed pages with a handwritten signature page is handled correctly.

\tikzfig[0.82]{activity_ingestion}{Activity Diagram -- Document Ingestion}{fig:act-ingest}

\subsubsection{Module 4: Answering a Question}
Figure~\ref{fig:act-query} shows the three answer paths. All paths start from the user's visible
documents.

\tikzfig[0.8]{activity_query}{Activity Diagram -- Answering a Question}{fig:act-query}

\subsubsection{Module 3: Review and Verification}
\tikzfig[0.6]{activity_review}{Activity Diagram -- Review and Verification}{fig:act-review}

\subsubsection{Module 1: Organization and Member Management}
\tikzfig[0.7]{activity_users}{Activity Diagram -- Organization and Member Management}{fig:act-users}

\subsection{Class Diagram}
Figure~\ref{fig:class} shows the main classes, their key attributes and operations, and their
relationships. Filled diamonds mark composition: versions, pages, chunks and fields cannot exist without
their parent.

\tikzfig[0.8]{class_diagram}{Class Diagram of \sysname{}}{fig:class}

\subsection{Sequence Diagrams}
\subsubsection{Sequence Diagram -- User Login}
\tikzfig[0.5]{seq_login}{Sequence Diagram -- User Login}{fig:seq-login}

\subsubsection{Sequence Diagram -- Core Operation (Asking a Question)}
Figure~\ref{fig:seq-ask} shows a question passing through the router, the access filter and one of
the answer engines. Only the LOOKUP and KNOWLEDGE alternatives are drawn; SEARCH follows the KNOWLEDGE
path without the final LLM call.

\tikzfig[0.62]{seq_ask}{Sequence Diagram -- Asking a Question}{fig:seq-ask}

\subsection{State Transition Diagrams}
A document version moves through the states in Figure~\ref{fig:state-version}. Because a new version
becomes current only when it is READY, the previous version stays searchable while the new one is
processed. An extracted field moves through the states in Figure~\ref{fig:state-field}; every
transition caused by a person creates a FieldEdit record.

\tikzfig[0.4]{state_version}{State Transition Diagram -- Document Version Lifecycle}{fig:state-version}
\tikzfig[0.4]{state_field}{State Transition Diagram -- Record Field Lifecycle}{fig:state-field}

\section{Data Flow Diagrams}
The context diagram (Figure~\ref{fig:context}) is the Level 0 DFD. The Level 1 DFD
(Figure~\ref{fig:dfd1}) divides the system into five processes and seven data stores, and the Level 2
DFD (Figure~\ref{fig:dfd2}) expands process 4.0, answering questions.

\subsection{Level 1 Data Flow Diagram}
\tikzfig[0.6]{dfd1}{Level 1 Data Flow Diagram}{fig:dfd1}

\subsection{Level 2 Data Flow Diagram}
\tikzfig[0.5]{dfd2}{Level 2 Data Flow Diagram -- Process 4.0 Answer Questions}{fig:dfd2}

\section{Data Design}
\subsection{Entity Relationship Diagram (ERD)}
\tikzfig[0.75]{erd}{Entity Relationship Diagram of \sysname{}}{fig:erd}

\subsection{Database Architecture}
The database is normalized to Third Normal Form. Key design decisions are:
\begin{itemize}
\item \textbf{Document and version are separate.} A \emph{Document} is the logical item (``NDA with
Acme''); each uploaded file is a \emph{DocumentVersion}. Pages, chunks and fields belong to a version,
so history is preserved exactly and versions can be compared. \code{current\_version} points to the
newest READY version.
\item \textbf{Typed field values.} \code{RecordField} keeps the value as text and, depending on the field
type, also as a \code{date} or \code{decimal}, so questions such as ``expiring in 2027'' are answered
with indexed range queries instead of text matching.
\item \textbf{Search structures inside PostgreSQL.} The \code{embedding} column (\code{vector(384)})
has an HNSW index (\code{m=16, ef\_construction=64}, cosine), and \code{search\_vector} (section path
weighted A, text weighted B) has a GIN index.
\item \textbf{Integrity rules in the database.} Unique constraints exist on (document, number),
(version, key), (organization, department name) and (recipient, reminder key). The audit model refuses
updates and deletes.
\end{itemize}

\section{Database Schema Diagram}
\tikzfig[0.85]{schema}{Database Schema Diagram of \sysname{}}{fig:schema}

\section{Data Dictionary}
The data dictionary for the most important tables is given below. Django adds an auto-increment
primary key \code{id} to every table.

\begin{table}[H]\centering\small
\caption{Data Dictionary -- DOCUMENT Table}\label{tab:dd-doc}
\begin{tabularx}{\textwidth}{|p{3.2cm}|p{2.6cm}|X|}\hline
\textbf{Column} & \textbf{Type} & \textbf{Description} \\ \hline
organization\_id & FK, not null & Owning organization \\ \hline
department\_id & FK, nullable & Department the document belongs to \\ \hline
title & varchar(255) & Display title \\ \hline
kind & varchar(10) & UNKNOWN, RECORD or REFERENCE \\ \hline
kind\_locked & boolean & True when a user set the kind; not overwritten by the classifier \\ \hline
doc\_type & varchar(120) & e.g.\ ``Non-Disclosure Agreement'', ``Law / statute'' \\ \hline
summary & text & One-sentence summary \\ \hline
access\_level & varchar(10) & ORG, DEPARTMENT or RESTRICTED \\ \hline
uploaded\_by\_id & FK User & Uploader (always allowed to see it) \\ \hline
current\_version\_id & FK Version & Newest READY version \\ \hline
is\_archived & boolean & Archived documents are hidden everywhere \\ \hline
use\_as\_policy & boolean & Reference document whose rules records are checked against \\ \hline
created\_at, updated\_at & timestamp & Audit timestamps \\ \hline
\end{tabularx}
\end{table}

\begin{table}[H]\centering\small
\caption{Data Dictionary -- DOCUMENT VERSION Table}\label{tab:dd-ver}
\begin{tabularx}{\textwidth}{|p{3.2cm}|p{2.6cm}|X|}\hline
\textbf{Column} & \textbf{Type} & \textbf{Description} \\ \hline
document\_id & FK, not null & Parent document \\ \hline
number & integer & 1, 2, 3 \dots; unique per document \\ \hline
file, original\_filename & varchar & Stored path and original name \\ \hline
file\_type & varchar(5) & PDF, DOCX or IMAGE \\ \hline
file\_size, sha256 & integer, char(64) & Size in bytes and content hash \\ \hline
page\_count & integer & Number of pages read \\ \hline
marked\_handwritten & boolean & Uploader says the document is handwritten \\ \hline
change\_note & varchar(300) & Reason for the new version \\ \hline
status & varchar(10) & QUEUED, PROCESSING, READY, FAILED \\ \hline
processing\_log & JSON & Per-step timings and details (method, page sources, counts) \\ \hline
change\_summary & text & AI summary of what changed from an earlier version \\ \hline
change\_summary\_from & integer & Number of the version it was compared with \\ \hline
policy\_checked\_at & timestamp & When the version was last checked against company policy \\ \hline
\end{tabularx}
\end{table}

\begin{table}[H]\centering\small
\caption{Data Dictionary -- POLICY FINDING Table}\label{tab:dd-policy}
\begin{tabularx}{\textwidth}{|p{3.2cm}|p{2.6cm}|X|}\hline
\textbf{Column} & \textbf{Type} & \textbf{Description} \\ \hline
version\_id & FK, not null & Record version that was checked \\ \hline
rule & varchar(300) & Short name of the rule, e.g.\ ``Notice periods'' \\ \hline
status & varchar(10) & OK, ISSUE or UNCLEAR \\ \hline
detail & text & One sentence comparing the record's fact with the rule \\ \hline
policy\_document\_id & FK Document & Policy document that holds the rule \\ \hline
policy\_page & integer & Page of the rule in that document \\ \hline
policy\_quote & text & The rule's words, found in the policy text (grounded) \\ \hline
created\_at & timestamp & When the check ran \\ \hline
\end{tabularx}
\end{table}

\begin{table}[H]\centering\small
\caption{Data Dictionary -- CHUNK Table}\label{tab:dd-chunk}
\begin{tabularx}{\textwidth}{|p{3.2cm}|p{2.6cm}|X|}\hline
\textbf{Column} & \textbf{Type} & \textbf{Description} \\ \hline
version\_id & FK & Version the text came from \\ \hline
sequence & integer & Order inside the document \\ \hline
section\_path & varchar(500) & e.g.\ ``Chapter II $>$ Section 25 $>$ Agreement without consideration, void'' \\ \hline
heading & varchar(300) & Heading of the chunk's section \\ \hline
text & text & Chunk text (at most 1,500 characters) \\ \hline
page\_start, page\_end & integer & Pages covered (used for citations) \\ \hline
clause\_label, clause\_confidence & varchar, float & Clause type from the fine-tuned model \\ \hline
embedding & vector(384) & MiniLM sentence embedding, HNSW-indexed \\ \hline
search\_vector & tsvector & Weighted full-text vector, GIN-indexed \\ \hline
\end{tabularx}
\end{table}

\begin{table}[H]\centering\small
\caption{Data Dictionary -- RECORD FIELD Table}\label{tab:dd-field}
\begin{tabularx}{\textwidth}{|p{3.2cm}|p{2.6cm}|X|}\hline
\textbf{Column} & \textbf{Type} & \textbf{Description} \\ \hline
version\_id & FK & Version the value was extracted from \\ \hline
key, label & varchar & One of the 11 field keys, and its display label \\ \hline
value\_text & text & Value as shown to users \\ \hline
value\_date & date & Parsed date for date fields \\ \hline
value\_number & decimal(18,2) & Parsed amount for money fields \\ \hline
confidence & float & 0--1; halved when the quote is not found on the page \\ \hline
source\_page, source\_quote & integer, text & Evidence location and verbatim quote \\ \hline
grounded & boolean & Quote found in the page text \\ \hline
status & varchar(12) & AUTO, NEEDS\_REVIEW, VERIFIED, CORRECTED \\ \hline
\end{tabularx}
\end{table}

\begin{table}[H]\centering\small
\caption{Data Dictionary -- AUDIT EVENT Table}\label{tab:dd-audit}
\begin{tabularx}{\textwidth}{|p{3.2cm}|p{2.6cm}|X|}\hline
\textbf{Column} & \textbf{Type} & \textbf{Description} \\ \hline
organization\_id & FK & Organization the event belongs to \\ \hline
actor\_id, actor\_name & FK, varchar & Who acted; the name is kept even if the user is deleted \\ \hline
action & varchar(20) & LOGIN, UPLOAD, NEW\_VERSION, VIEW, DOWNLOAD, SEARCH, FIELD\_EDIT, FIELD\_VERIFY, UPDATE\_DOCUMENT, ACCESS\_CHANGE, ARCHIVE, REPROCESS, MEMBER\_ADD, MEMBER\_UPDATE, DEPARTMENT\_CHANGE \\ \hline
target\_type, target\_id, target\_label & varchar & What the action was applied to \\ \hline
metadata & JSON & Extra details (e.g.\ old and new value) \\ \hline
ip\_address & inet & Client address \\ \hline
created\_at & timestamp, indexed & Time of the event \\ \hline
\end{tabularx}
\end{table}

\begin{table}[H]\centering\small
\caption{Data Dictionary -- QUERY Table}\label{tab:dd-query}
\begin{tabularx}{\textwidth}{|p{3.2cm}|p{2.6cm}|X|}\hline
\textbf{Column} & \textbf{Type} & \textbf{Description} \\ \hline
conversation\_id & FK & Conversation (owned by one user) \\ \hline
text & text & The question \\ \hline
intent & varchar(10) & LOOKUP, KNOWLEDGE, SEARCH or COMPARE \\ \hline
plan & JSON & Router output: filters and method (LLM or heuristic) \\ \hline
scope\_document\_ids & JSON & Documents the user limited the question to \\ \hline
answer, sources, results & text, JSON, JSON & Answer, cited passages/fields, record rows or hits \\ \hline
latency\_ms & integer & Response time \\ \hline
\end{tabularx}
\end{table}

\section{Database Tables}
\subsection{User and Organization Tables}
Django's \code{auth\_user} table stores credentials (PBKDF2 password hashes). The
\code{accounts\_membership} table links each user to exactly one organization (unique
\code{user\_id}), one role and optionally one department. All capability checks
(\code{can\_upload}, \code{can\_verify}, \code{can\_view\_audit}, \code{sees\_everything}) are
derived from the role in one place.

\subsection{Core Data Tables}
The core tables are \code{document}, \code{documentversion}, \code{page}, \code{chunk},
\code{recordfield} and \code{fieldedit}, described above, together with \code{processingjob} (the job
queue), \code{conversation} and \code{query} (query history), \code{auditevent} and
\code{notification}.

\subsection{Relationships Between Tables}
An organization has many departments, members, documents, conversations, audit events and
notifications. A document has many versions; a version has many pages, chunks, fields and jobs; a field
has many edits. Documents and users have a many-to-many ``shared with'' relationship for restricted
documents. A conversation has many queries.

\subsection{Security and Data Integrity}
\begin{itemize}
\item Every query on documents is restricted by organization and access rules in
\code{visible\_documents}; tests check the full role $\times$ access-level matrix.
\item Cascading deletes remove pages, chunks and fields with their version; audit events use
\code{SET\_NULL} for the actor and keep \code{actor\_name}, so history survives user deletion.
\item The audit model raises an error on update or delete.
\item Files are stored under non-guessable names and are only returned by an endpoint that checks access
and writes a DOWNLOAD or VIEW event.
\end{itemize}

\subsection{Scalability Considerations}
HNSW gives sub-linear vector search; the GIN index does the same for keyword search. Queries use
\code{hnsw.iterative\_scan} so that permission filtering does not reduce the number of results. The
worker can be scaled to several processes because jobs are claimed with \code{SKIP LOCKED}, and the API
is stateless (JWT), so more API containers can be added behind the web server.

\section{Chapter Summary}
This chapter presented the design of \sysname{}: the design methodology and process model, the layered
architecture with a background worker, activity, class, sequence and state diagrams, data flow diagrams,
the ERD, the schema and the data dictionary. The next chapter describes how this design was implemented.
